\documentclass[11pt]{article}
\setlength{\textwidth}{6.3in}
\setlength{\oddsidemargin}{0.1in}
\usepackage{amsmath,amssymb,amsthm}
\usepackage{hyperref}
\newtheorem{target}{Proposed theorem}
\newtheorem{lemma}{Lemma}
\newcommand{\Pcal}{\mathcal P}
\newcommand{\E}{\mathbb E}
\newcommand{\Aconst}{\mathfrak A}
\title{Squarefree Differences Between Primes:\\Spectral Gap and Parity Rigidity}
\author{Research draft}
\date{November 9, 2026}
\begin{document}
\maketitle
\begin{abstract}
We investigate the graph on the odd primes up to \(X\) in which two
vertices are joined when their difference is squarefree. The proposed
result identifies the constant and modulo-four parity functions as its
two dominant modes, with a sharp residual factor \(1/5\). Its intended
consequences are a sharp density bound and a quantitative inverse
theorem for prime sets avoiding squarefree differences.
\end{abstract}

\noindent\textbf{Proof status.}
The main spectral theorem is a proposed result under audit.
The uniform finite tail bound, concrete finite unit model, tuple occupancy
and cardinality, modulo-four sign adapter, and the complete finite
actual-prime comparison~(7g) have kernel-checked proofs in the maintained
Lean repository. The latter is uniform over every finite odd-prime
coordinate set, endpoint and signed real weight, with all exception terms.
The fixed-modulus occupancy limit, complete asymptotic tail passage,
sharpness and parity stability remain to be formalized.

\section{Statement and intended consequences}

Let \(\Pcal_X=\{p\le X:p\text{ is an odd prime}\}\) and
\(n_X=|\Pcal_X|\). All averages below are uniform averages on this
finite set; statements concern sufficiently large \(X\), so \(n_X>0\).
Write
\[
 \sigma(p)=
 \begin{cases}1,&p\equiv1\pmod4,\\-1,&p\equiv3\pmod4,\end{cases}
 \qquad
 \Aconst=\prod_{\ell\ {\rm prime}}
       \left(1-\frac1{\ell(\ell-1)}\right)>0 .
\]
For an arbitrary real function \(f\) on \(\Pcal_X\), define
\[
 a=\E f,\qquad b=\E(\sigma f),\qquad e=\E f^2,\qquad
 Q_X(f)=\frac1{n_X^2}\sum_{p,q\in\Pcal_X}
   f(p)f(q)\,\mathbf1_{\{|p-q|\ {\rm squarefree}\}}.
\]
Zero is not squarefree, so the diagonal contribution is zero.

\begin{target}[Sharp signed spectral estimate]\label{target:spectrum}
For every \(\varepsilon>0\) there is \(X_0\) such that for every
integer \(X\ge X_0\) and every \(f:\Pcal_X\to\mathbb R\),
\[
 \left|Q_X(f)-\Aconst(a^2-b^2)\right|
 \le \frac{\Aconst}{5}(e-a^2-b^2)+\varepsilon e.       \tag{1}
\]
The constant \(1/5\) is optimal. The threshold \(X_0\) is independent of \(f\).
\end{target}

The finite-\(X\) parity classes need not have exactly equal sizes.
Thus (1), including its error term, is the proposed statement;
no exact finite-\(X\) orthogonality of \(1\) and \(\sigma\) is assumed.

For a set \(B_X\subseteq\Pcal_X\) with no squarefree difference between
distinct elements, put
\[
 \alpha_X=|B_X|/n_X,\qquad
 \beta_X=\E(\sigma\mathbf1_{B_X}).
\]
Substitution of \(f=\mathbf1_{B_X}\) into (1) would give, uniformly in \(B_X\),
\[
 6\alpha_X^2-4\beta_X^2\le\alpha_X+o(1).              \tag{2}
\]
Since \(|\beta_X|\le\min(\alpha_X,1-\alpha_X)+o(1)\), this implies
\(\alpha_X\le1/2+o(1)\). Either reduced class modulo four supplies
a matching construction: every difference within that class is divisible by four.
For \(1/6<\alpha_X\le1/2\), (2) further gives
\[
 \frac{\min_{r\in\{1,3\}}|B_X\mathbin\triangle
       \{p\in\Pcal_X:p\equiv r\pmod4\}|}{n_X}
 \le \frac12-\frac12\sqrt{6\alpha_X^2-\alpha_X}+o(1). \tag{3}
\]
In particular, a density deficit \(\delta\) from \(1/2\) gives an
asymptotic symmetric-difference bound
\(5\delta/4+O(\delta^2)+o(1)\).

\section{Existing analytic input}

The maintained programme supplies one positive constant \(C\), chosen
before the modulus, character and endpoint, such that
\[
 \left|\sum_{p\le X}\chi(p)\right|
 \le C X^{7/8}\bigl(1+\log m+\log(X+2)\bigr)^2        \tag{4}
\]
for every nonprincipal Dirichlet character modulo \(m\ge1\).
This is the already verified seven-eighths input and its prime-prefix
consequence. The principal character is retained separately.
The exact maintained owner is
\texttt{Helpers/QRHPrimePrefix.lean}, declaration
\texttt{exists\_primeCharacterSum\_le\_seven\_eighths}.
Character orthogonality then bounds the count in any reduced residue class
by \(\pi(X)/\varphi(m)\) plus the error in (4).
This use includes the nonunit correction in the principal center.

\section{The first arithmetic test: a uniform tail}

For a prime \(p\le X\) and integers \(2\le R\le T\), let
\[
 H_{R,X}(p)=\#\{q\le X:q\text{ prime},\ q\ne p,\
        \exists\text{ prime }\ell>R:\ell^2\mid(p-q)\}.
\]
The required estimate is uniform in the moving vertex \(p\):
\[
 H_{R,X}(p)\le
 \frac{\pi(X)}R+
 C T X^{7/8}\bigl(1+\log(T^2)+\log(X+2)\bigr)^2+
 2\left\lfloor\frac XT\right\rfloor .                \tag{5}
\]
Here the constant \(C\) is independent of \(p,R,T,X\).

\paragraph{Proof of the finite bound (5).}
For \(R<\ell\le T\), if \(\ell\ne p\), the residue of \(p\) modulo
\(\ell^2\) is a unit, so (4) and residue reconstruction apply.
If \(\ell=p\), the off-diagonal count is zero: divisibility of \(p-q\)
by \(p^2\) would imply \(p\mid q\), hence \(q=p\).
The full low-square union is bounded by summing its exact counts.
Its principal contribution is at most
\[
 \pi(X)\sum_{R<\ell\le T}\frac1{\ell(\ell-1)}
 \le \pi(X)\sum_{j=R+1}^{T}\frac1{j(j-1)}
 \le \frac{\pi(X)}R .
\]
There are at most \(T\) low moduli, giving the middle term of (5).
For the high-square union, split at \(q=p\) and use
\(0<|p-q|\le X\). Each sign has at most
\[
 \sum_{u>T}\left\lfloor\frac X{u^2}\right\rfloor
 \le\left\lfloor\frac XT\right\rfloor
\]
possible positive differences. This retains every high divisor and
does not add a rounding error for each divisor.
With \(T=\max(2,\lfloor X^{1/16}\rfloor)\), (5) yields
\[
 \limsup_{X\to\infty}\ \max_{\substack{p\le X\\p\ {\rm prime}}}
       \frac{H_{R,X}(p)}{\pi(X)}\le\frac1R .         \tag{6}
\]
The actual limit argument also uses the ordinary prime number theorem.
It does not replace the moving vertex by a fixed prime.

\section{Finite spectral model and proposed prime transfer}

\begin{lemma}[Complete-product gap with the binary mode retained]\label{lem:finite}
Let \(m_1,\ldots,m_k\ge6\), allowing \(k=0\), and put
\(\Omega=\{0,1\}\times\prod_{j=1}^k\{1,\ldots,m_j\}\).
Give \(\Omega\) uniform measure and let
\[
 K((b,x),(c,y))=\mathbf1_{\{b\ne c\}}
                   \prod_{j=1}^k\mathbf1_{\{x_j\ne y_j\}},\qquad
 D=\frac12\prod_{j=1}^k(1-m_j^{-1}).
\]
For every real \(f\) on \(\Omega\), put
\(a=\E f\), \(b=\E((-1)^{x_0}f)\) and \(e=\E f^2\). Then
\[
 \left|\E_{x,y} f(x)K(x,y)f(y)-D(a^2-b^2)\right|
       \le \frac D5(e-a^2-b^2).
\]
\end{lemma}

\begin{proof}
First omit the binary coordinate and normalize each complete-graph
operator by its degree, so its rows sum to one. On a support of size
\(N\), write
\[
 B(u,v)=\sum_x u(x)(Pv)(x),\qquad
 S(u)=\sum_x u(x),\qquad E(u)=\sum_x u(x)^2.
\]
We prove by induction over all the coordinates that
\[
 |B(u,u)-S(u)^2/N|\le\tfrac15(E(u)-S(u)^2/N).       \tag{7a}
\]
The one-vertex base has zero variance. For a new coordinate of size
\(m\ge6\), write its slices as \(f_i=g+r_i\), where
\(g=m^{-1}\sum_i f_i\) and \(\sum_i r_i=0\) pointwise.
Bilinearity, with both cross terms retained, gives exactly
\[
 \begin{split}
 \sum_{i\ne j}B(f_i,f_j)&=m(m-1)B(g,g)-\sum_iB(r_i,r_i),\\
 \sum_i E(f_i)&=mE(g)+\sum_i E(r_i),\qquad
 \sum_iS(f_i)=mS(g).
 \end{split}                                                   \tag{7b}
\]
The full double sum is \(m^2B(g,g)\); its diagonal is
\(mB(g,g)+\sum_iB(r_i,r_i)\). Thus no row or overlap has been discarded.

Equation (7a) and nonnegative centered energy imply
\(|B(v,v)|\le E(v)\) for every \(v\).
After division of the first line of (7b) by \(m-1\), the centered
new form is the sum of
\[
 m\bigl(B(g,g)-S(g)^2/N\bigr)
       \quad\hbox{and}\quad -\frac1{m-1}\sum_iB(r_i,r_i).
\]
Their bounds are respectively \(m/5\) times the centered energy of \(g\)
and at most \(1/5\) times \(\sum_i E(r_i)\), because \(m-1\ge5\).
The energy and mean identities in (7b) prove (7a) for the enlarged product.

For the binary coordinate, write its slices as \(f_0=g+h\) and
\(f_1=g-h\). Its form is the signed expression
\[
 B(f_0,f_1)+B(f_1,f_0)=2B(g,g)-2B(h,h).
\]
Apply (7a) to both \(g\) and \(h\).
The energy identity is \(E(f_0)+E(f_1)=2E(g)+2E(h)\);
the two means are \(S(f_0)+S(f_1)=2S(g)\) and
\(S(f_0)-S(f_1)=2S(h)\).
Dividing by the total number \(2N\) of vertices yields the two-mode
bound for the row-normalized operator. Multiplication by its uniform
row mass \(D\) gives the displayed inequality.
The Lean declarations check both the row-normalized induction and the
zero-one kernel with its density factor. Explicitly, for
\(N=\prod_j m_j\) and \(d=\prod_j(1-m_j^{-1})\), the raw odd-product
kernel is \(Nd\) times its row-normalized kernel.
The binary extension has \(2N\) vertices and density \(D=d/2\).
Its checked counting form is exactly
\[
 \begin{aligned}
 &\left|\sum_{x,y\in\Omega}f(x)K(x,y)f(y)
       -D\left[(\sum_{x\in\Omega}f(x))^2
              -(\sum_{x\in\Omega}(-1)^{x_0}f(x))^2\right]\right|\\
 &\qquad\le\frac D5\left[
 |\Omega|\sum_{x\in\Omega}f(x)^2
 -(\sum_{x\in\Omega}f(x))^2
 -(\sum_{x\in\Omega}(-1)^{x_0}f(x))^2\right].
 \end{aligned}
\]
Dividing this inequality by \(|\Omega|^2\) gives the displayed
probability form. The empty odd product has \(N=d=1\), so the
binary-only case is included without a separate convention.
\end{proof}

\subsection{Concrete prime-square unit model}

\begin{lemma}[Exact unit kernel and parity coordinate]\label{lem:unit-model}
Let \(J\) be any finite index set and let \(\ell_j\ge3\) be primes.
Put
\[
 \Omega_J=\{0,1\}\times
     \prod_{j\in J}(\mathbb Z/\ell_j^2\mathbb Z)^\times,\qquad
 D_J=\frac12\prod_{j\in J}
       \left(1-\frac1{\ell_j(\ell_j-1)}\right).
\]
The primes may repeat; \(J=\varnothing\) is included.
For \(u=(c,x)\), \(v=(d,y)\), define
\[
 L_J(u,v)=\mathbf1_{\{c\ne d\}}
            \mathbf1_{\{\forall j\in J,\ x_j\ne y_j\}},
 \qquad \eta(c,x)=(-1)^c .
\]
Then \(0<D_J\le1/2\), and every real \(h\) on \(\Omega_J\)
satisfies
\[
 \left|\E_{u,v}h(u)L_J(u,v)h(v)
        -D_J\bigl((\E h)^2-(\E(\eta h))^2\bigr)\right|
 \le \frac{D_J}{5}
       \bigl(\E h^2-(\E h)^2-(\E(\eta h))^2\bigr).
                                                               \tag{7f}
\]
The unit group modulo four is explicitly equivalent to \(\{0,1\}\):
the units \(1,3\) correspond respectively to \(0,1\).
Under this equivalence, unit equality is residue equality and
the sign \(+1,-1\) is exactly \((-1)^c\).
\end{lemma}
\begin{proof}
For each coordinate the unit group has exactly
\(\ell_j(\ell_j-1)\ge6\) elements.
Equality of units is equivalent to equality of their underlying residues.
Lemma~\ref{lem:finite} therefore applies directly to these dependent
coordinate types and gives (7f), with the displayed density.
Every odd prime contributes a factor strictly between zero and one;
the empty product is one, proving \(0<D_J\le1/2\).
Modulo four the only units are \(1\) and \(-1=3\).
Sending these to \(0\) and \(1\) gives an explicit bijection, with inverse
\(c\mapsto1\) for \(c=0\) and \(c\mapsto3\) for \(c=1\).
The equality, sign and inverse-residue assertions follow in these two cases.
The dependent-product estimate, its arithmetic specialization and each
of these three modulo-four identities have separate checked declarations.
\end{proof}

\subsection{Uniform finite-fibre transfer, with all exceptions}

\begin{lemma}[Signed sampling comparison]\label{lem:fibre}
Let \(S\) have size \(n>0\), let \(T\subseteq S\), and let \(U\) have
size \(M>0\). Let \(g:T\to U\), and suppose a real model kernel \(L\)
with sign function \(\eta\) satisfies
\[
 \begin{aligned}
 &\left|\sum_{u,v}h(u)L(u,v)h(v)
       -D\left[(\sum_u h(u))^2-(\sum_u\eta(u)h(u))^2\right]\right|\\
 &\qquad\le\frac D5\left[M\sum_u h(u)^2
       -(\sum_u h(u))^2-(\sum_u\eta(u)h(u))^2\right].
 \end{aligned}
\]
for every real \(h\), where \(D\ge0\).
Let \(K:S\times S\to\mathbb R\) satisfy \(|K|\le1\), with
\(K(p,q)=L(g(p),g(q))\) on \(T\times T\).
Let \(\sigma(p)^2=1\) on \(S\), with
\(\sigma(p)=\eta(g(p))\) on \(T\).
Assume \(\delta,\tau\ge0\) and
\[
 M\,|\{p\in T:g(p)=u\}|\le(1+\delta)n\quad(u\in U),
 \qquad |S\setminus T|\le\tau^2n.
\]
For uniform averages on \(S\), put \(a=\E f\), \(b=\E(\sigma f)\),
\(e=\E f^2\) and
\(q=n^{-2}\sum_{p,q\in S}f(p)K(p,q)f(q)\). Then, uniformly for every real \(f\),
\[
 |q-D(a^2-b^2)|
 \le \frac D5(e-a^2-b^2)
       +\left[\frac{D\delta}{5}+(4D+2)\tau+\tau^2\right]e.       \tag{7c}
\]
\end{lemma}

\begin{proof}
For \(u\in U\), set \(F(u)=\sum_{p\in T:g(p)=u}f(p)\) and
\(h(u)=MF(u)/n\). Empty fibres have total zero.
The signed regular quadratic form is exactly
\[
 q_T=\frac1{M^2}\sum_{u,v}h(u)L(u,v)h(v).
\]
Cauchy--Schwarz inside each fibre, followed by the occupancy bound, gives
\[
 \frac1M\sum_u h(u)^2
 =\frac M{n^2}\sum_u F(u)^2
 \le\frac{1+\delta}{n}\sum_{p\in T}f(p)^2\le(1+\delta)e.
\]
The two model means are exactly
\(a_T=n^{-1}\sum_{p\in T}f(p)\) and
\(b_T=n^{-1}\sum_{p\in T}\sigma(p)f(p)\).
Thus the model comparison gives
\[
 |q_T-D(a_T^2-b_T^2)|
       \le\frac D5((1+\delta)e-a_T^2-b_T^2).
\]
Write \(E_0=S\setminus T\). The exact remaining kernel packet consists of
\(T\times E_0\), \(E_0\times T\) and \(E_0\times E_0\).
Its absolute contribution is at most \((2\tau+\tau^2)e\);
each mixed rectangle is bounded by \(\tau e\) and the last by \(\tau^2e\).
No portion of the signed regular kernel is bounded this way.
The same sparse-support estimate gives
\[
 |\Delta_a|\le2\tau e,\qquad|\Delta_b|\le2\tau e,\qquad
 \Delta_a=a^2-a_T^2,\quad\Delta_b=b^2-b_T^2.
\]
Here the parity estimate applies the ordinary squared-mean bound to
\(\sigma f\), whose energy equals that of \(f\).

Subtract the desired residual main term before bounding these shifts.
For the upper inequality their coefficients are
\(-4D/5\) on \(\Delta_a\) and \(6D/5\) on \(\Delta_b\).
For the lower inequality the two coefficients interchange.
Consequently their combined allowance is \(4D\tau e\).
Adding the complete exceptional packet proves (7c).
The proof is formalized in
\texttt{Finset.kernel\_fibre\_fifth\_gap\_with\_exceptions};
its two finite mean corrections and all exceptional rectangles are retained.
\end{proof}

\subsection{Residue model and pending prime transfer}

Fix \(R\ge3\). Put
\[
 U_R=\prod_{\ell\le R}(\mathbb Z/\ell^2\mathbb Z)^\times,\quad
 r_\ell=\varphi(\ell^2)=\ell(\ell-1),\quad
 D_R=\prod_{\ell\le R}(1-r_\ell^{-1}).
\]
Equip \(U_R\) with uniform product measure. Its kernel is
\(K_R(u,v)=\prod_{\ell\le R}\mathbf1_{\{u_\ell\ne v_\ell\}}\).
At coordinate \(\ell\), the averaging operator has eigenvalue
\(1-r_\ell^{-1}\) on constants and \(-r_\ell^{-1}\) on the
mean-zero space. Orthogonal tensor decomposition therefore gives eigenvalues
\[
 D_R\prod_{\ell\in S}\frac{-1}{r_\ell-1}
       \qquad (S\subseteq\{\ell:\ell\le R\}).        \tag{7}
\]
The two modes from \(S=\varnothing,\{2\}\) are \(1,\sigma\)
with eigenvalues \(D_R,-D_R\).
Every other mode has modulus at most \(D_R/5\), since
\(r_\ell-1\ge5\) for odd \(\ell\). Consequently the finite kernel satisfies
the counterpart of (1) with no error and \(D_R\) in place of \(\Aconst\).

\begin{lemma}[Uniform upper occupancy of prime fibres]\label{lem:prime-fibres}
There is \(C>0\) such that the following holds for every integer modulus
\(m\ge1\), every nonnegative integer \(X\), every finite subset \(T\) of the primes
at most \(X\) coprime to \(m\), and every map \(g:T\to U\) satisfying
\[
 g(p)=g(q)\quad\Longrightarrow\quad p\equiv q\pmod m.
\]
For every \(u\in U\),
\[
 |\{p\in T:g(p)=u\}|\le\frac{\pi(X)}{\varphi(m)}
 +C X^{7/8}(1+\log m+\log(X+2))^2.                 \tag{7d}
\]
The constant is independent of \(m,X,T,U,g,u\); no inverse or
surjectivity of \(g\) is assumed.
\end{lemma}
\begin{proof}
For a nonempty fibre choose a member \(p_0\). Every other member
lies in the reduced residue class of \(p_0\) modulo \(m\).
The independent, uniform reduced-residue bound following (4) therefore
bounds the fibre by the right side of (7d), with its same constant.
For an empty fibre the assertion follows from nonnegativity of that
right side. This also covers modulus one.
\end{proof}

\begin{lemma}[Actual prime-square tuples]\label{lem:prime-square-tuples}
Let \(S\) be any finite set of primes and put
\[
 N_S=\prod_{\ell\in S}\ell^2,\qquad
 M_S=\prod_{\ell\in S}\ell(\ell-1),\qquad
 T_S(X)=\{q\le X:q\text{ prime},\ q\notin S\}.
\]
Then \(\varphi(N_S)=M_S\), and the full unit-tuple space
\(\prod_{\ell\in S}(\mathbb Z/\ell^2\mathbb Z)^\times\)
has cardinality \(M_S\). Every \(q\in T_S(X)\) is a unit modulo
\(N_S\). Equality of all prime-square residue coordinates implies
congruence modulo \(N_S\).

There is one constant \(C>0\), before \(S\), the nonnegative integer
\(X\) and the residue tuple \(u\), such that
\[
 \#\{q\in T_S(X):(q\bmod\ell^2)=u_\ell\ (\ell\in S)\}
 \le\frac{\pi(X)}{M_S}
 +C X^{7/8}(1+\log N_S+\log(X+2))^2.              \tag{7e}
\]
Here \(u\) ranges over the full residue-tuple space, so no inverse
choice or surjectivity hypothesis is imposed.
\end{lemma}
\begin{proof}
Distinct coordinate primes have coprime squares. Induction over \(S\)
and the injectivity of the two-factor Chinese remainder map therefore
combine equal coordinate residues into one congruence modulo \(N_S\).
The same coprimality and
\(\varphi(\ell^2)=\ell(\ell-1)\) give the totient formula.
The cardinality of a dependent product of finite unit groups is the
product of their cardinalities, giving \(M_S\).
A prime outside \(S\) is coprime to every coordinate square and hence
to \(N_S\). Lemma~\ref{lem:prime-fibres} now applies to the actual tuple
map, giving (7e) with its unchanged constant.

All assertions include \(S=\varnothing\): then \(N_S=M_S=1\)
and there is one empty tuple. The coordinate \(\ell=2\) contributes
\(\varphi(4)=2\); its prime is excluded from \(T_S(X)\) exactly when
\(2\in S\), as are all other coordinate primes.
\end{proof}

\begin{lemma}[Complete finite comparison on actual primes]\label{lem:prime-finite}
There is one \(C>0\) with the following property.
Let \(S\) be any finite set of odd primes, and put
\[
 N_S=4\prod_{\ell\in S}\ell^2,\qquad
 M_S=2\prod_{\ell\in S}\ell(\ell-1),\qquad
 D_S=\frac12\prod_{\ell\in S}
       \left(1-\frac1{\ell(\ell-1)}\right).
\]
For an integer \(X\ge0\) with \(n=n_X>0\), define
\[
 E_{S,X}=CX^{7/8}(1+\log N_S+\log(X+2))^2,\qquad
 \delta_{S,X}=\frac{1+M_SE_{S,X}}n .
\]
On \(\Pcal_X\), let
\[
 K_S(p,q)=\mathbf1_{\{p\not\equiv q\pmod4\}}
     \prod_{\ell\in S}\mathbf1_{\{p\not\equiv q\pmod{\ell^2}\}},\qquad
 q_S(f)=\frac1{n^2}\sum_{p,q\in\Pcal_X}f(p)K_S(p,q)f(q).
\]
For every \(\tau\ge0\) with \(|S|\le\tau^2n\), and every real
\(f\) on \(\Pcal_X\), the original means \(a,b,e\) satisfy
\[
 |q_S(f)-D_S(a^2-b^2)|
 \le\frac{D_S}{5}(e-a^2-b^2)
  +\left[\frac{D_S\delta_{S,X}}5+
           (4D_S+2)\tau+\tau^2\right]e.                    \tag{7g}
\]
The constant precedes \(S,X,\tau,f\).
The empty set \(S\) is included, with \(N_S=4,M_S=2,D_S=1/2\).
\end{lemma}
\begin{proof}
Take \(T=\Pcal_X\setminus S\).
Every member is a unit modulo four and every \(\ell^2\), \(\ell\in S\).
Map it to the binary coordinate supplied by the modulo-four equivalence
and its odd prime-square unit coordinates.
Lemma~\ref{lem:unit-model} identifies both the physical kernel on
\(T\times T\) and the sign \(\sigma\) with that model.

Apply (7e) to \(\{2\}\cup S\).
Equality of the unit labels implies equality of all those actual residue
coordinates. Thus each label fibre has size at most
\(\pi(X)/M_S+E_{S,X}\); empty fibres have size zero.
The model has \(M_S\) elements, and
\[
 M_S|T_u|\le\pi(X)+M_SE_{S,X}
 \le n+1+M_SE_{S,X}=(1+\delta_{S,X})n .
\]
The \(1/n\) term accounts for using the full-prime center \(\pi(X)\)
with the odd-prime vertex count \(n\).
There are at most \(|S|\le\tau^2n\) exceptional vertices, and
\(|K_S|\le1\) everywhere.
All hypotheses of (7c) therefore hold on the same prime support and with
the same real weight \(f\). Its complete exception allowance gives (7g).
In particular the two mixed rectangles and the exceptional self
interaction, as well as both signed principal means, are retained.
\end{proof}

For fixed \(R\), put \(Q_R=\prod_{\ell\le R}\ell^2\) and
\(M_R=\varphi(Q_R)=|U_R|\).
Let \(T\) consist of the odd primes at most \(X\) that do not divide \(Q_R\).
Lemma~\ref{lem:fibre} requires only an upper bound on each fibre.
For a nonempty fibre choose one prime \(p_0\) in it. Every other member
satisfies \(p\equiv p_0\pmod{\ell^2}\) for all \(\ell\le R\);
pairwise coprimality gives \(p\equiv p_0\pmod{Q_R}\).
The representative is a unit modulo \(Q_R\).
An empty fibre has count zero.
This avoids constructing a global Chinese-remainder inverse merely to count fibres.

The actual tuple bound (7e) gives
\[
 |T_u|\le\frac{\pi(X)}{M_R}
    +C X^{7/8}(1+\log Q_R+\log(X+2))^2.
\]
Since \(n_X=\pi(X)-1\) for \(X\ge2\), the occupancy condition holds with
\[
 \delta_{R,X}=\frac1{n_X}
   +\frac{M_R C X^{7/8}(1+\log Q_R+\log(X+2))^2}{n_X}
       \longrightarrow0
\]
as \(X\to\infty\) with \(R\) fixed.
The exceptional count is at most the number of odd primes dividing \(Q_R\).
For every fixed \(\tau>0\), it is eventually at most \(\tau^2 n_X\).
Equation (7g) supplies the complete actual-prime finite-kernel comparison,
uniformly in \(f\), with these exact occupancy and exception errors.
The arithmetic tuple assertions, unit model and parity adapter are
checked separately, and their complete actual-prime finite assembly is
now checked as Lemma~\ref{lem:prime-finite}.
The fixed-\(R\) limit of its error remains pending Lean work;
the actual-prime asymptotic headline theorem remains proposed.

Let \(J_R\) be the difference between this finite kernel and the actual
squarefree-difference kernel. It is a symmetric nonnegative tail kernel,
and (5) bounds every row sum. Only this separated tail is estimated by
\[
 \left|\frac1{n_X^2}\sum_{p,q}J_R(p,q)f(p)f(q)\right|
 \le \left(\max_p\frac1{n_X}\sum_qJ_R(p,q)\right)\E f^2,
\]
using \(2|f(p)f(q)|\le f(p)^2+f(q)^2\).
All signs of the principal finite kernel remain in (7).
First let \(X\) tend to infinity with \(R\) fixed, then let \(R\)
tend to infinity. Equation (6) removes the tail and
\(D_R\to\Aconst\), giving the proposed conclusion (1).

For sharpness, take a nonzero real mean-zero function of the unit residue
modulo nine alone, with arbitrary bounded values at the exceptional prime three.
Its constant and parity averages tend to zero. Its local eigenvalue at three
is \(-1/6\), so (7) gives the limiting Rayleigh quotient
\(-\Aconst/5\).

\section{Relation to existing work and verification plan}

Finite direct-product independence and stability are classical inputs;
see Alon--Defant~\cite{AD}. Arithmetic graph limits based on integer
GCD data are studied by Martineau~\cite{M}.
The proposed contribution here is the sharp signed operator transfer
to the squarefree-difference graph on actual primes and its uniform
prime-subset inverse consequence.

The finite bound is formalized as
\texttt{exists\_uniform\_prime\_square\_tail\_upper} in the namespace
\texttt{PrimeFactorOscillations.QRH}, with maintained owner
\texttt{Helpers/QRHPrimeSquareTail.lean}.
The formal count uses \(q\equiv p\pmod{\ell^2}\), the congruence formulation
of the divisibility condition in~(5). The source compiles through the public
export; the eight supporting theorems use only the standard foundational axioms.
The finite model now has maintained proofs in
\texttt{CompleteCoordinateGap.lean} and \texttt{CompleteProductGap.lean},
under \texttt{Mathlib/Analysis/MeanInequalities}.
The declarations \texttt{Finset.IsLargeCompleteProduct.fifth\_gap} and
\texttt{Finset.IsLargeCompleteProduct.parity\_gap} prove the arbitrary-arity
centered gap and the separate binary mode.
The construction predicate carries only finite graph-product data and
coordinate cardinality conditions; it assumes no spectral estimate.
The full eigenvalue description (7) is an ordinary spectral explanation,
not an additional Lean classification claim.
The finite-fibre theorem~\ref{lem:fibre}, including its full exceptional
packet, is now checked in \texttt{FiberKernelGap.lean}, supported by
\texttt{KernelRestriction.lean} in the same candidate directory.
The zero-one density conversion is checked in
\texttt{RawProductSpectrum.lean}: the declarations
\texttt{Finset.IsRawCompleteProduct.normalized\_witness},
\texttt{zero\_or\_one}, \texttt{parity\_gap} and
\texttt{binary\_model\_gap} retain the full density and both modes.
Lemma~\ref{lem:prime-fibres} is checked as
\texttt{PrimeFactorOscillations.QRH.exists\_prime\_fibre\_count\_upper}
in \texttt{Helpers/QRHPrimeFibreOccupancy.lean}.
The arithmetic assertions of Lemma~\ref{lem:prime-square-tuples} are
checked in \texttt{Mathlib/Data/Nat/Prime/SquareResidues.lean}, by
\texttt{Nat.coprime\_prime\_sq\_prod},
\texttt{cast\_zmod\_prime\_sq\_prod\_eq},
\texttt{totient\_prod\_prime\_sq},
\texttt{isUnit\_cast\_prod\_prime\_sq} and
\texttt{card\_prime\_sq\_unit\_tuples}.
The actual tuple count (7e) is checked as
\texttt{PrimeFactorOscillations.QRH.}
\texttt{exists\_prime\_square\_tuple\_fibre\_count\_upper}
in \texttt{Helpers/QRHPrimeSquareFibres.lean}.
These are separate checked declarations; the combined lemma above is
their mathematical consequence.
The concrete dependent-product estimate is checked in
\texttt{FinitePiProduct.lean}, by
\texttt{Finset.completePiDensity\_bounds} and
\texttt{kernelBilinear\_complete\_pi\_binary\_gap}.
Lemma~\ref{lem:unit-model} has arithmetic declarations
\texttt{Nat.prime\_square\_unit\_density\_bounds} and
\texttt{prime\_square\_unit\_binary\_gap} in
\texttt{Mathlib/Data/Nat/Prime/SquareUnitKernel.lean}.
The same owner checks
\texttt{ZMod.unitsFourEquivBool\_eq\_iff},
\texttt{unitsFourSign\_eq\_binary} and
\texttt{unitsFourEquivBool\_symm\_val}.
All seven new declarations compile and their focused axiom audit contains
only the standard foundational axioms.
These finite model statements are supporting classical mathematics,
rather than the paper's claimed originality.
The complete finite actual-prime comparison~(7g) is checked as
\texttt{PrimeFactorOscillations.QRH.}
\texttt{exists\_finite\_prime\_square\_spectral\_gap}, in
\texttt{Helpers/QRHPrimeSquareSpectrum.lean}.
Its focused axiom audit contains only the standard foundational axioms.
The theorem includes every finite odd-prime coordinate set, the empty
coordinate set, arbitrary real weights, the full-prime center correction,
and every coordinate-prime exception.
Remaining formalization consists of the fixed-modulus occupancy limit,
the complete asymptotic tail passage, sharpness and consequences (2)--(3).
This document is kept separate from the metaplectic and zero-multiplicity papers.

\begin{thebibliography}{9}
\bibitem{AD} N. Alon and C. Defant,
\emph{Isoperimetry, Stability, and Irredundance in Direct Products},
Discrete Mathematics (2020).
\url{https://web.math.princeton.edu/~nalon/PDFS/isostab.pdf}.
\bibitem{M} S. Martineau,
\emph{On coprime percolation, the visibility graphon, and the local limit
of the GCD profile}, arXiv:1804.06486, revised 2019.
\url{https://arxiv.org/abs/1804.06486}.
\end{thebibliography}
\end{document}
